\documentclass{article}
\usepackage[UTF8]{ctex}
\usepackage{amsmath,amssymb,amsfonts}
\usepackage{enumitem}
\usepackage{geometry}
\geometry{a4paper, margin=1in}
% === 错题本排版优化（不参与题目标识） ===
\setlength{\parindent}{0pt}
\setlength{\parskip}{0.35em plus 0.10em minus 0.08em}
\setlength{\emergencystretch}{3em}
\setlist{itemsep=0.25em, topsep=0.40em, parsep=0pt}
\allowdisplaybreaks[2]
% === 错题本排版优化结束 ===
\begin{document}
\section*{南京理工大学2025年高等代数真题}
\section*{90．南京理工大学2025年高等代数试卷}
一．填空题．

\begin{itemize}

\subsection{题目1}
  \item[1．] 已知
\[
D=\left|\begin{array}{ccc}
0 & -1 & 2 \\
2 & 5 & 8 \\
0 & 5 & 3
\end{array}\right|=-26 .
\]
则 \(2 A_{31}-5 A_{32}+8 A_{33}=\) \(\_\_\_\_\)，其中 \(A_{i j}\) 是元素 \(a_{i j}\) 的代数余子式．

\subsection{题目2}
  \item[2．] 若 5 阶矩阵 \(A\) 的行列式 \(|A|=-1, A^{*}\) 是 \(A\) 的伴随矩阵，则 \(\left|A^{*}\right|=\) \(\_\_\_\_\)．
\end{itemize}

\subsection{题目3}
3．已知线性方程组 \(\left(\begin{array}{ccc}1 & 2 & 1 \\ 2 & 3 & a+2 \\ 1 & a & -2\end{array}\right)\left(\begin{array}{l}x_{1} \\ x_{2} \\ x_{3}\end{array}\right)=\left(\begin{array}{l}1 \\ 3 \\ 0\end{array}\right)\) 无解，则 \(a=\)

\subsection{题目4}
4．设 \(A\) 为2阶矩阵，\(\alpha_{1}, \alpha_{2}\) 为线性无关的2维列向量，\(A \alpha_{1}=0, A \alpha_{2}=2 \alpha_{1}+\alpha_{2}\) ，则 \(A\) 的非零特征值为 \(\_\_\_\_\)

\subsection{题目5}
5．秩为3的5元二次型只有 \(\_\_\_\_\)中规模型。\\
\(\_\_\_\_\)

\subsection{题目6}
6．已知实二次型

\[
f\left(x_{1}, x_{2}, x_{3}\right)=a x_{1}^{2}+x_{2}^{2}+4 a x_{3}^{2}+6 x_{1} x_{2}+5 x_{1} x_{3}-8 x_{2} x_{3}
\]

经正交变换可化为标准型 \(f=-16 y_{1}^{2}\) ，则 \(a=\) \(\_\_\_\_\)．

\subsection{题目7}
7．设 \(A\) 是 \(m \times n\) 矩阵，则齐次线性方程组 \(A X=0\) 仅有零解得充要条件是 \(\_\_\_\_\)．

\subsection{题目8}
8．\(f(x)=4 x^{4}+2 x^{3}+2 x^{2}+x, g(x)=2 x^{3}-4 x^{2}+x-2\) ，则 \(f(x)\) 与 \(g(x)\) 的首一最大公因式 =\\
\(\_\_\_\_\)．

二．设 \(A, B\) 均为2阶矩阵，\(A^{*}, B^{*}\) 分别为 \(A, B\) 的伴随矩阵，若 \(|A|=-2,|B|=-4\) ，请给出分块矩阵

\[
\left(\begin{array}{ll}
O & A \\
B & O
\end{array}\right)
\]

的伴随矩阵（用 \(A^{*}, B^{*}\) 表示）。

三．已知平面上三条不同直线的方程分别为

\[
\begin{aligned}
& l_{1}: a x+b y+c=0 \\
& l_{2}: b x+c y+a=0 \\
& l_{3}: c x+a y+b=0 .
\end{aligned}
\]

试证三条直线相交于一点的充分必要条件是 \(a+b+c=0\) ．

四．设3阶实对称矩阵 \(A\) 的各行元素之和均为2，向量 \(\alpha_{1}=(1,0,-1)^{\mathrm{T}}, \alpha_{2}=(2,-1,-1)^{\mathrm{T}}\) 是线性方程组 \(A X=0\) 的两个解。

\begin{itemize}
  \item[（1）] 求 \(A\) 的特征值和特征向量．
  \item[（2）] 求正交矩阵 \(Q\) 和对角矩阵 \(\Lambda\) ，使得 \(Q^{\mathrm{T}} A Q=\Lambda\) 。
\end{itemize}

五．求 \(\mathbb{R}^{3}\) 的基 \(\alpha_{1}=(1,-1,0), \alpha_{2}=(0,1,0), \alpha_{3}=(1,0,-1)\) 的对偶基 \(f_{1}, f_{2}, f_{3}\) ．

六．设 \(A, B\) 是同阶正交阵，\(|A|=-|B|\) ，证明：\(A+B\) 是奇异矩阵。

七．设 \(V\) 为 \(n\) 维线性空间，证明：\(V\) 的任一非平凡子空间都是若干个 \(n-1\) 维子空间的交。

八．设 \(f(x)\) 是整系数多项式，它在 9 个整数点处取值均为 8，证明：\(f(x)\) 无整数根．
\end{document}
